% =========================================================
% Continuity — Notes
% Chapter: continuity
% Section: continuity
% Volume III · Analysis
% =========================================================

\subsection{Continuity}
\label{sec:continuity}

% ---------------------------------------------------------
% Toolkit
% ---------------------------------------------------------

\begin{tcolorbox}[title={Toolkit --- Continuity}]
\begin{itemize}
  \item Three equivalent forms: $\varepsilon$-$\delta$, neighbourhood,
    sequential. The $\varepsilon$-$\delta$ form is unpunctured
    ($x=c$ allowed), unlike function limits.
  \item $\delta$ may depend on both $\varepsilon$ and $c$.
    If $\delta$ depends only on $\varepsilon$, the condition is
    uniform continuity.
  \item Isolated points are automatically continuous.
  \item Algebra: sums, differences, products, scalar multiples,
    and quotients (denominator nonzero) of continuous functions
    are continuous.
  \item Composition preserves continuity.
\end{itemize}
\end{tcolorbox}

% ---------------------------------------------------------
% Definition — Relative Neighborhood
% ---------------------------------------------------------

\begin{tcolorbox}[title={Definition --- Relative Neighborhood}]
\begin{definition}[Relative Neighborhood]
\label{def:relative-neighborhood}
Let $E \subseteq \mathbb{R}$ and let $a \in E$. A
\textbf{neighborhood of $a$ in $E$} is a set of the form
\[
U_E(a) := E \cap (a-\delta, a+\delta)
\quad \text{for some } \delta > 0.
\]
\end{definition}
\end{tcolorbox}

\begin{remark*}[Standard quantified statement]
\[
\begin{aligned}
&U_E(a) \text{ is a neighborhood of } a \text{ in } E \\
&\quad\Longleftrightarrow\quad
\exists \delta > 0 :\;
U_E(a) = E \cap (a-\delta, a+\delta).
\end{aligned}
\]
\end{remark*}

\begin{remark*}[Interpretation]
A relative neighborhood restricts an open interval in $\mathbb{R}$
to the set $E$.
\end{remark*}

% ---------------------------------------------------------
% Definition — Continuity at a Point
% ---------------------------------------------------------

\begin{tcolorbox}[title={Definition --- Continuity at a Point}]
\begin{definition}[Continuity at a Point]
\label{def:continuous-at-point}
Let $A \subseteq \mathbb{R}$, let $f : A \to \mathbb{R}$, and let
$c \in A$. The function $f$ is \textbf{continuous at $c$} if
\[
\forall \varepsilon > 0 \;\exists \delta > 0 \;\forall x \in A :
\bigl(|x-c| < \delta \Rightarrow |f(x)-f(c)| < \varepsilon\bigr).
\]
\end{definition}
\end{tcolorbox}

\begin{remark*}[Standard quantified statement]
\[
\begin{aligned}
&f \text{ is continuous at } c \\
&\quad\Longleftrightarrow\quad
\forall \varepsilon > 0 \;\exists \delta > 0 \;\forall x \in A : \\
&\qquad (|x-c| < \delta \Rightarrow |f(x)-f(c)| < \varepsilon).
\end{aligned}
\]
\end{remark*}

\begin{remark*}[Interpretation]
Continuity means small changes in input produce small changes in output.
\end{remark*}

% ---------------------------------------------------------
% Definition — Continuity (Neighbourhood)
% ---------------------------------------------------------

\begin{tcolorbox}[colback=defbox, colframe=defborder, arc=2pt,
  left=6pt, right=6pt, top=4pt, bottom=4pt,
  title={\small\textbf{Definition (Continuity at a Point --- Neighbourhood Form)}},
  fonttitle=\small\bfseries]
\begin{definition}[Continuity at a Point --- Neighbourhood Form]
\label{def:continuous-at-point-nbhd}
Let $A \subseteq \mathbb{R}$, let $f : A \to \mathbb{R}$, and let
$c \in A$. The function $f$ is continuous at $c$ if
\[
\forall \varepsilon > 0 \;\exists \delta > 0 :
f(A \cap (c-\delta, c+\delta))
\subseteq (f(c)-\varepsilon, f(c)+\varepsilon).
\]
\end{definition}
\end{tcolorbox}


\begin{remark*}[Standard quantified statement]
\[
\begin{aligned}
&f \text{ is continuous at } c \text{ (nbhd)} \\
&\quad\Longleftrightarrow\quad
\forall \varepsilon > 0 \;\exists \delta > 0 \;\forall x \in A : \\
&\qquad (x \in (c-\delta, c+\delta)
\Rightarrow f(x) \in (f(c)-\varepsilon, f(c)+\varepsilon)).
\end{aligned}
\]
\end{remark*}

\begin{remark*}[Interpretation]
The neighborhood around $c$ maps into a neighborhood of $f(c)$.
\end{remark*}

% ---------------------------------------------------------
% Definition — Continuity (Sequential)
% ---------------------------------------------------------

\begin{tcolorbox}[colback=defbox, colframe=defborder, arc=2pt,
  left=6pt, right=6pt, top=4pt, bottom=4pt,
  title={\small\textbf{Definition (Continuity at a Point --- Sequential Form)}},
  fonttitle=\small\bfseries]
\begin{definition}[Continuity at a Point --- Sequential Form]
\label{def:continuous-at-point-seq}
Let $A \subseteq \mathbb{R}$ and $c \in A$. The function $f$ is
continuous at $c$ if
\[
\forall (x_n) \subseteq A :
x_n \to c \Rightarrow f(x_n) \to f(c).
\]
\end{definition}
\end{tcolorbox}


\begin{remark*}[Standard quantified statement]
\[
\begin{aligned}
&f \text{ is continuous at } c \text{ (seq)} \\
&\quad\Longleftrightarrow\quad
\forall (x_n) \subseteq A :
x_n \to c \Rightarrow f(x_n) \to f(c).
\end{aligned}
\]
\end{remark*}

\begin{remark*}[Interpretation]
Continuity preserves limits of sequences.
\end{remark*}
